\scanpage{9}
% source lines 1-4
Or on a $\operatorname{div}(\vec y)\equiv\operatorname{div}(a)$ (car $\operatorname{div}(\vec x)=\unsure{1|0}$) et il faut donc prendre les $\operatorname{div}(a)$, où $a\in K^*$ est tel que $a_{\mathfrak p}\in V_{\mathfrak p,\underline m}$ si $\mathfrak p\mid\underline m$, --- \unsure{ce qui s'écrit}
\[
a\equiv1\ (\underline m).
\]

On a par suite

\deleted{$\unsure{C\simeq J/V_{\underline m}K^*}$}
\[
C_{\underline m}(K)=C(K)/H_{\underline m}(K)\simeq\unsure{J|\mathscr{J}}(K)/V_{\underline m}(K)K^*\simeq{\unsure{D|\mathscr{D}}}_{\underline m}(K)/P_{\underline m}(K)
\]

% source lines 7-10
avec

$P_{\underline m}(K)=$ \unsure{classes de diviseurs} de la forme $(a)$, avec $a\equiv1\ (\underline m)$ (ce qui \deleted{signifie} \inserted{implique} que $a$ est \uline{positif} aux places à l'infini réelles).

$H_{\underline m}(K)=V_{\underline m}(K)K^*/K^*$. On a aussi

% source: exact sequence and its two annotations
\diagram{\begin{tikzcd}[ampersand replacement=\&,column sep=small]
0\arrow[r]\&\substack{\text{ss-groupe de }\deleted{\unsure{E(K)^*}}\inserted{E(K)}\text{ formé}\\\text{des éléments qui sont}\\\equiv1\ (\underline m)\\\text{[ce qui implique}\\\text{qu'ils sont }\geq0\text{ aux places}\\\text{réelles]}}\arrow[r]\&K^*\arrow[r]\&\deleted{\unsure{P}}\underbrace{D_{\underline m}(K)\times\displaystyle\prod_{\substack{\deleted{\unsure{\mathfrak p\text{ fini}}}\\\mathfrak p\mid\underline m}}}\underbrace{K_{\mathfrak p}/V_{\mathfrak p,\underline m}}_{\substack{\mathbb Z/2\text{ aux places}\\\text{réelles, }0\text{ aux}\\\text{places complexes,}\\\text{sinon une extension}\\\text{de }\mathbb Z\text{ par un groupe}\\\deleted{\text{\unsure{d'ordre}}\ \unsure{Z_{\mathfrak p}}}\ ......}}\arrow[r]\&C_{\underline m}(K)\arrow[r]\&0
\end{tikzcd}}

% source lines 11-15
Soit $x\in K_{\mathfrak p}$, où $\mathfrak p$ est une place quelconque.

Alors $x$ définit un \unsure{élmt|èlmt|élément} de $J(K)$, donc de $C(K)$ donc de $C_{\underline m}(K)$, donc de $D_{\underline m}(K)/P_{\underline m}(K)$.

Pour le calculer, \deleted{\unsure{on sait que}} \deleted{\unsure{on fait la chose}} \unsure{il vient} $\operatorname{div}(x)\bmod P_{\underline m}(K)$ si $\mathfrak p\nmid\underline m$. \uline{Si} $\mathfrak p\mid\underline m$, $\longrightarrow$
% archive sheet number: 8
